\documentclass[10pt,a4paper]{amsart}
\usepackage[margin=19mm]{geometry}
\usepackage{amsmath,amssymb,mathtools}
\usepackage[scr=rsfs,cal=euler]{mathalfa}
\usepackage{xfrac}
\usepackage[unicode,hidelinks,bookmarks=false]{hyperref}
\newcommand{\ie}{i.e.\ }
\newcommand{\cf}{cf.\ }
\newcommand{\resp}{respectively\ }
\newcommand{\Subsection}{\subsection}
\providecommand{\nobreakdash}{\nobreak\hbox{-}}
\setlength{\emergencystretch}{2em}
\multlinegap0pt
\usepackage{amssymb}
\usepackage{float}
\newtheorem{proposition}{Proposition}
\newcommand{\Adjointaction}{\mathrm{Ad}}
\newcommand{\Bxi}{\mathcal{B}_{\xi}}
\newcommand{\Cinf}{C^\infty}
\newcommand{\Hamiltonian}[1]{\mathcal{H}_{#1}}
\newcommand{\Reducedxi}{\mathcal{R}_\xi}
\newcommand{\SET}[1]{\{#1\}}
\newcommand{\galgebra}{\mathfrak{g}}
\newcommand{\gstar}{\mathfrak{g}^*}
\title{Superintegrability of stratified symplectic spaces}
\author{Zhuo Chen, Kai Jiang$^*$, Nicolai Reshetikhin, Husileng Xiao}
\date{Source: arXiv:2609.01953v1 (CC BY 4.0)}
\begin{document}
\maketitle

\noindent\emph{Source and licence.} Superintegrability of stratified symplectic spaces. Version arXiv:2609.01953v1 (CC BY 4.0), distributed by arXiv under the Creative Commons Attribution 4.0 International licence, http://creativecommons.org/licenses/by/4.0/, as recorded in the arXiv licence metadata for this exact version and shown on the abstract page https://arxiv.org/abs/2609.01953v1. Full licence notice: \texttt{licenses/arxiv-2609.01953-CC-BY-4.0.txt}.\par\smallskip

\noindent\textit{Editorial scope.} Counted unit: the article's proposition Prop:tangenttorefinedstrata (source0.tex lines 1159--1162). The article's complete original proof of it is delivered with it (source0.tex lines 1164--1206, one \texttt{proof} environment of 43 lines). No mathematics is written, repaired or re-derived: the statement and the proof are the article's own lines, in the article's own notation, and no retained line is substituted or re-flowed.  No further source-local context is needed: every cross-reference of every retained line resolves inside the two delivered spans.  Nothing was omitted from any retained span, so the package carries no omission record.  No retained line contains a citation, so the wrapper carries no bibliography and no bibliography entry.  No retained line contains a figure environment, an image-inclusion command or drawing code, so no image is delivered and the package declares no asset dependency.  Editorial formatting only, in addition: an AMS \texttt{amsart} preamble that restates the article's own declarations and macros used by the retained lines, namely the environments \texttt{proposition} on separate counters, together with the standard packages those lines need.  The retained environments print in this document as Proposition 1 in order of appearance, whereas the article prints the same items with the numbering of its own full document.  The complete native source of the article as distributed in the arXiv e-print for this exact version is bundled unchanged as \texttt{sources/209155/source0.tex}.
\begin{proposition}\label{Prop:tangenttorefinedstrata}
  The Hamiltonian vector field $\Hamiltonian{h}$ generated by $h \in \Cinf(\Bxi)$ is tangent to every refined stratum $(\Reducedxi)_{(H_1),\beta,(H_2)}$.

\end{proposition} 

\begin{proof}
  Since we have identified $T^*G$ with $ \gstar\times G$ (via right translations), the algebra   $C^\infty(T^*G)$ is generated by two types of basic functions:  (1) vectors $X\in \galgebra$, seen as linear functions on $\gstar$ and on $T^*G$ as well; (2) functions $f\in C^\infty(G)$, seen as functions on $T^*G$ which do not change along the $\gstar$-direction. As for the Poisson bracket relations, $\SET{X_1,X_2}$ coincides with $[X_1,X_2]_{\galgebra}$, $\SET{f_1,f_2}$ is simply zero, and $\SET{X,f}$ is $\overrightarrow{X}(f)$. Therefore, the Hamiltonian vector field $\Hamiltonian{X}$ of $X\in \galgebra$ on $T^*G\cong \gstar \times G$ is expressed as follows: \begin{equation}\label{Eqt:HamiltonianXsingle}
\Hamiltonian{X}|_{(x,\gamma)}=(-\mathrm{ad}_X^*(x),\overrightarrow{X}|_{\gamma}),\quad \forall (x,\gamma)\in \gstar\times G.            
         \end{equation}
Here we have identified $ T_{(x,\gamma)}(\gstar\times G)$ with $\gstar\times T_{\gamma}G$.
 One could also consider 
 $$X=\sum_i X_{i_1}\odot X_{i_2} \odot \cdots \odot X_{i_n}\in S^n(\galgebra), \quad(\mbox{the $n$-th symmetric product of }\galgebra)  $$
where every $X_{\cdot}$ is a vector in $\galgebra$. This $X$ is also regarded as a function on $\gstar$, and on $\gstar\times G$ as well. Based on \eqref{Eqt:HamiltonianXsingle}, we can inductively show that
\begin{equation}\label{Eqt:HamiltonianXmulti}
\Hamiltonian{X}|_{(x,\gamma)}=
(-\mathrm{ad}_{x \llcorner X}^*(x),\overrightarrow{x\llcorner X}|_{\gamma}).            
         \end{equation}
         Here, we define 
         $$x\llcorner X:=\sum_i \sum_{j=1}^n \langle X_{i_1},x \rangle \langle X_{i_2},x \rangle   \cdots \widehat{X_{i_j}} \cdots \langle X_{i_n}, x \rangle X_{i_j} \in \galgebra. $$

When $X\in S^n(\galgebra)$ is $G$-invariant, we can regard $X$ as a function on $\gstar/G$, and also on $\Bxi$.  However, $X$ being $G$-invariant implies that 
\begin{equation}\label{Eqt:tempXGinvairantzero}
 [X,Y]:=\sum_i \sum_{j=1}^n  X_{i_1}  \odot X_{i_2}    \cdots \widehat{X_{i_j}} \cdots \odot X_{i_n} \odot [X_{i_j},Y]=0,\quad \forall Y \in \galgebra,
\end{equation}
and hence 
$\mathrm{ad}_{x \llcorner X}^*(x)=0$. Therefore, according to \eqref{Eqt:HamiltonianXmulti}, only the second component,  i.e., the $T_xG$-direction, of $\Hamiltonian{X}|_{(x,\gamma)}$  is nontrivial. The flow generated by $\Hamiltonian{X}$ is given by
\begin{equation}\label{Eqt:flowX} \varphi^t_X(x,\gamma)=(x,(\exp(t ~x\llcorner X))\gamma).
\end{equation}
 In this expression, the element $\exp(t ~x\llcorner X)$ belongs to $G_x$ because we can show the following identity:
 \begin{equation}\label{Eqt:temphere}
  \langle Y, \Adjointaction^*_{\exp(t ~x\llcorner X)}(x) \rangle \equiv \langle Y, x \rangle,\quad \forall Y\in \galgebra, t\in \mathbb{R}.
 \end{equation}
 
 
  In fact, the left-hand side of Equation \eqref{Eqt:temphere} can be alternatively computed by:
  \begin{eqnarray*}
    \langle  \Adjointaction_{\exp(-t ~x\llcorner X)}(Y),x \rangle &=& \langle  \exp( -t ~\mathrm{ad}_{x\llcorner X} ) Y,x \rangle \\
      &=& \langle Y, x \rangle+\sum_{j=1}^\infty \frac{(-t)^j}{j!} \langle \mathrm{ad}^j_{x\llcorner X}Y, x \rangle.
  \end{eqnarray*}
  It can be inductively shown that
  \begin{equation}\label{Eqt:tempadtoj}
   \langle \mathrm{ad}^j_{x\llcorner X}Y, x \rangle= {[X,\mathrm{ad}^{j-1}_{x\llcorner X}Y ]} (x),\qquad j\geqslant 1.
   \end{equation}
  Here the right-hand side is interpreted as the value of the function $[X,\mathrm{ad}^{j-1}_{x\llcorner X}Y ]\in S^n{(\galgebra)}\subset C^\infty(\gstar)$ on $x\in \gstar$. However, since we have assumed the $G$-invariant property of $X$, namely Equation \eqref{Eqt:tempXGinvairantzero}, the above Equation \eqref{Eqt:tempadtoj} simply yields zero, and \eqref{Eqt:temphere} is justified.
  
  Finally, from Equation \eqref{Eqt:flowX}  and $\exp(t ~x\llcorner X) \in G_x$, it can be easily seen that   $[(x,\gamma)]\in (\Reducedxi)_{(H_1),\beta,(H_2)}$ implies $\varphi^t_X[(x,\gamma)]=[\varphi^t_X(x,\gamma)]\in (\Reducedxi)_{(H_1),\beta,(H_2)}$ as well. Since every function $h \in \Cinf(\Bxi)$ can be (locally) approximated by  polynomial functions on $ \gstar$, the flow $\varphi^t_h$ satisfies the same property.
   
\end{proof}

\end{document}
